\section{Lecture 4 (September 16th)}
\begin{thm}
Let $D$ be a domain in ${\bm C}$. There is a sequence $\{K_{n}\}$ of compact subsets of $D$ and $\{\delta _{n}\}$ ($\delta _{n}>0$) such that
\begin{itemize}
\item[(i)] \[\bigcup ^{\infty }_{n}K_{n}=D\]
\item[(ii)] If $K\subset D$ is compact then $K\subset K_{n}$ for some $n$
\item[(iii)] $\bar{D}(z,2\delta _{n})\subset K_{n+1}$ for every $z\in K_{n}$, $n=1,2,\ldots $
\end{itemize}
Note that $z\notin F$ and $F$ is closed. This implies that $d(z,F)=\min\{|z-w| \,|\, w\in F\}>0$.
\end{thm}
\vspace{2ex}
\begin{proof}
Put
\[K_{n}=\Big\{z\leq n\Big\}\cap \Big\{d(z,D^{C})\geq \dfrac{1}{n}\Big\}\]
\end{proof}
\vspace{2ex}
\begin{ex}
Consider 
\[f_{n}(x)=\dfrac{\sin (nx)}{\sqrt{n}}\rightarrow 0\]
which converges uniformly on ${\bm R}$. However,
\[f'(0)=\sqrt{n}\rightarrow \infty \]
as $n\rightarrow \infty $.
\end{ex}
\vspace{2ex}
\begin{thm}
If $\{f_{n}\}$ is a sequence of analytic functions on a domain $D$ such that $f_{n}\rightarrow f$ uniformly on every compact subset of $D$, then 
\begin{itemize}
\item[(i)] $f$ is analytic on $D$ (by Morera's theorem)
\item[(ii)] $f_{n}'\rightarrow f'$ uniformly on every compact subset of $D$
\end{itemize}
\end{thm}
\vspace{2ex}
\begin{proof}
Take any compact subset $K$ of $D$. Then $K\subset K_{n}$ for some $n$, and $\bar{D}(z,\delta _{n})\subset K_{n+1}$ for all $z\in K$. So, there exists $r>0$ and $K'$, a compact subset of $D$, such that $\bar{D}(z,r)\subset K'$ for all $z\in K$. Here, applying Cauchy's integral formula,
\[f(z)=\dfrac{1}{2\pi i}\int _{|\xi -z|=r}\dfrac{f(\xi )}{(\xi -z)^2}\,d\xi \]
For every $z\in K$, we have
\[f'(z)=\dfrac{1}{2\pi i}\int _{|\xi -z|=r}\dfrac{f(\xi )}{(\xi -z)^2}\,d\xi \]
which implies that
\[|f'(z)|\leq \dfrac{1}{2\pi }\dfrac{||f||_{K'}}{r^2}(2\pi r)=\dfrac{||f||_{K'}}{r}\]
with $||f||_{K}=\max_{z\in K}|f(z)|$. Therefore, for every $z\in K$,
\[|f_{n}'(z)-f'(z)|\leq \dfrac{||f_{n}-f||_{K'}}{r}\]
wihch means $f_{n}'\rightarrow f'$ uniformly on $K$. 
\end{proof}
\vspace{2ex}
\begin{ex}
(Linear mapping) Consider the linear mapping
\[w=az+b\]
This mapping we see that, by expressing $a=re^{i\theta }$, the map stretches, rotates, and translates the complex plane.
\end{ex}
\vspace{2ex}
\begin{defi}
(Linear functional transformation) For $ad-bc\ne 0$,
\[w=T(z)=\dfrac{az+b}{cz+d}\]
This can be written as
\[w=T(z)=\dfrac{a}{c}+\dfrac{bc-ad}{c}\cdot \dfrac{1}{cz+d}\]
Additionally, the map is 1-1 from the fact that the following inverse map exists:
\[z=\dfrac{dw-b}{-cw+a}\]
Thus telling us that the map can be seen as a combination of translation, dilation ($|a|>1$) or contraction ($|a|\leq 1$), rotation, and inversion. 
\[z\rightarrow \dfrac{1}{z}\]
This is the tricky one! 
\end{defi}
\vspace{2ex}
\begin{rmk}
(Inversion map $w=1/z$) The equation (locus) of all lines and all circles in ${\bm R}^2$ is
\[a(x^2+y^2)+bx+cy+d=0\]
where $b^2+c^2>4ad $ and $a,b,c,d\in {\bm R}$. In the complex plane, for $x^2+y^2=|z|^2$, $x=(z+\bar{z})/2$ and $y=(z-\bar{z})/2i$ so that
\[a|z|^2+\dfrac{1}{2}(b-ic)z+\dfrac{1}{2}(b+ic)\bar{z}+d=0\]
and again $b^2+c^2>4ad$. This implies that
\[a|z|^2+\beta z+\bar{\beta }\bar{z}+d=0\]
where $a,d\in {\bm R}$ and $\beta \in {\bm C}$ with $|\beta |^2>ad$. Taking the inversion map,
\[a\dfrac{1}{|w|^2}+\beta\dfrac{1}{w}+\bar{\beta }\dfrac{1}{\bar{w}}+d|w|^2=0\]
and
\[a+\beta \bar{w}+\bar{\beta }w+d=0\]
by multiplying $|w|^2=w\bar{w}$ ($\beta \bar{\beta }>ad$). We therefore see that the inversion map takes the set of all lines and circles to again the set of all lines and circles.
\end{rmk}
\vspace{2ex}
\begin{defi}
($S^2$) $S^2$ is defined as the boundary of a unit ball in ${\bm R}^3$. We see this as a one-point compactification (Alexandroff compactification) of ${\bm C}$, and is seen as ${\bm C}\cup  \{\infty \}$ where $\{\infty \}$ is called the north pole. 
\end{defi}
\vspace{2ex}
\begin{lem}
The linear functional transformation which takes $\{a,b,c\}$ to $\{0,1,\infty \}$ is uniquely determined by
\[T(z)=\dfrac{(z-a)(b-c)}{(z-c)(b-a)}\]
\end{lem}
\vspace{2ex}

